% Study 1: how the number of topics was chosen.
% Written by tools/make_tables.py from the pipeline table(s) below (paths relative to
% Analysis/R Analysis Pipeline, each with the SHA-256 it was read with); do not edit by hand.
%   Study 1/tables/text_topics_stability.csv  fa312640420b5534bc1eceb367d6f094fbd1949ca6765c465870583d9c140475
\begin{table}[htbp]
\centering
\small
\caption[Study 1: how the number of topics was chosen]{\textbf{Study~1: how the number of topics was chosen.} For each number of topics from 3 to 10, how stable the topic model of each rating's explanations was (exploratory). The most stable number was used: 6 for both ratings (in bold).}\label{tab:supp_text_topics_stability}
\setlength{\tabcolsep}{4pt}
\begin{tabular}{@{}lll@{}}
\toprule
Number of topics & Creativity: stability (SD) & Effort: stability (SD)\\
\midrule
3 & 0.83 (0.13) & 0.86 (0.15)\\
4 & 0.89 (0.08) & 0.87 (0.14)\\
5 & 0.88 (0.07) & 0.94 (0.08)\\
6 & \textbf{0.92 (0.05)} & \textbf{0.98 (0.01)}\\
7 & 0.84 (0.07) & 0.95 (0.05)\\
8 & 0.82 (0.08) & 0.91 (0.06)\\
9 & 0.84 (0.06) & 0.87 (0.06)\\
10 & 0.87 (0.07) & 0.86 (0.06)\\
\bottomrule
\end{tabular}
\par\smallskip
\begin{minipage}{\textwidth}\footnotesize\raggedright A topic model (here non-negative matrix factorization) finds groups of words that tend to occur together. At each number of topics it was refitted ten times, each time on a random 80\% of the participants. The topics of two refits were matched one to one, and their agreement is the mean cosine similarity of the matched topics' word weights (1 when two topics weight the same words in the same proportions, 0 when they share no word). Stability is that agreement averaged over every pair of refits; SD, its standard deviation. A tie at three decimals goes to the smaller number.\par\end{minipage}
\end{table}
